\section{Version 1 and the audit}

Version 1 ingested 208 repositories found by three GitHub Search API queries (recently created and already
high-starred; old, moderate and recently maintained; not pushed in years), giving 414,271 daily rows from
2008-11-13 to 2026-06-22, of which 179 repositories had at least 270 days of history (41, 69 and 69 in the
three cohorts). The backtest used an expanding context starting 180 days after the first star, a 30-day step,
at most eight windows per repository, horizons of 7, 30 and 90 days, and 200 predictive trajectories per
model. Four classical baselines (horizon-dependent seasonal naive, linear trend, damped ETS, LightGBM) and two
foundation models, Chronos-2 and TimesFM-2.5, were scored by empirical CRPS. Both foundation models had lower
mean CRPS than every original baseline in all nine cohort-by-horizon cells; in the ``breakout'' cohort the
foundation models scored 6.2 to 8.5 against 11.2 to 28.0 for the baselines.

The audit of 2026-07-12 found no code defect that invalidated the aggregate six-model comparison; every
summary field recomputed from the saved artifact. It found two design flaws that invalidated the original
claims:
\begin{enumerate}
\item \emph{Outcome-conditioned cohorts.} ``Breakout'' was selected for already having many stars, ``dead''
  for later inactivity, ``steady'' for later maintenance. None of these labels was available at forecast
  time. The ``dead'' cohort's evaluation windows had a median of 5.02 stars per day and only 7.7 percent zero
  days; it was not dead when evaluated.
\item \emph{Post-peak windows.} The global peak preceded the evaluated span for 78.0 percent of breakout
  repositories; only 12.2 percent peaked inside it. The study tested the forecasting of post-viral decay, not
  of shocks.
\end{enumerate}
After the audit added a 28-day local mean on identical windows, the foundation models still had the lower
mean in every cell (Appendix Table~\ref{tab:v1}), but the per-repository win rate against it fell to between
40.6 and 80.5 percent, and the audit put the share of positive gain held by the five largest positive
contributors at roughly 30 to 77 percent depending on the cell, a figure it tied to no named reference. The
audit also listed what version 1 had to disclose (Appendix~\ref{app:v1}). Its prescription, freeze and
pre-register, context-only labels, trailing non-overlapping windows with a common endpoint, strong
pre-registered baselines, native quantile metrics, repository-clustered inference, and a prospective
post-release cohort as the contamination control, became the specification of version 2; the cohort is
registered and not yet run (Section~\ref{sec:notrun}). Every version 1 number in this report is quoted from
the preserved version 1 documents, not recomputed; the version 1 artifacts are on a computer the author no
longer uses.
